% Seminar 3 - Page 10. Fragment for
% % Seminar 3 - Page 10. Fragment for
% % Seminar 3 - Page 10. Fragment for
% % Seminar 3 - Page 10. Fragment for
% \input{seminars/pages/seminar3-pagina10-en}; no preamble or page break.
% EDITORIAL NOTES:
% - This final manuscript page continues the discussion of cardinality
%   begun on Seminar 3, Page 9.
% - For the cardinal 1 the handwritten set looks like {{empty set}}, rather
%   than the usual {empty set}. Either is a singleton and has cardinal 1.
% - The handwritten cardinal sum appears to use A_i UNION {i}, not
%   A_i TIMES {i}. That literal expression does not in general produce a
%   disjoint union, and its cardinality need not be the intended sum.
%   The standard tagged-disjoint-union construction would instead be
%     | \bigcup_{i\in I} (A_i \times \{i\}) |.
%   We retain the apparent handwriting in the displayed expression.
% - The manuscript states that these definitions are correct without a
%   proof of independence of the representing sets A_i, A and B.
% - "Cantor's paradox" is transcribed as a note about the putative set
%   of all cardinals; its proper-class distinction is not supplied in
%   the manuscript and is not inserted into the text.
% - The last comparison of cardinalities uses injective/bijective maps
%   without addressing any set-theoretic axioms. It is left as written.

\subsection*{Seminar 3 - Page 10}
\begingroup
\setlength{\abovedisplayskip}{5pt}
\setlength{\belowdisplayskip}{5pt}
\setlength{\abovedisplayshortskip}{3pt}
\setlength{\belowdisplayshortskip}{3pt}

\noindent The first cardinal numbers are written as
\[
\begin{aligned}
0 &:= \operatorname{card}(\varnothing),\\
1 &:= \operatorname{card}\bigl(\{\{\varnothing\}\}\bigr),\\
2 &:= \operatorname{card}\bigl(\{\{\varnothing\},\{\{\varnothing\}\}\}\bigr).
\end{aligned}
\]

\medskip
\noindent\underline{Operations with cardinal numbers.}

Let $\alpha=|A|$. If $(\alpha_i)_{i\in I}$ is a family of cardinal
numbers, with $\alpha_i=|A_i|$, define
\[
\sum_{i\in I}\alpha_i\ \overset{\mathrm{def}}{=}
\left|\bigcup_{i\in I}\bigl(A_i\cup\{i\}\bigr)\right|.
\]

\noindent For $\alpha=|A|$ and $\beta=|B|$, define
\[
\alpha\cdot\beta\ \overset{\mathrm{def}}{=}\ |A\times B|,
\qquad
\alpha^{\beta}\ \overset{\mathrm{def}}{=}
\bigl|\operatorname{Hom}(B,A)\bigr|.
\]

\noindent The definitions are well-defined.

\medskip
\noindent\underline{Cantor's paradox.}
The collection of all cardinal numbers, regarded as a set,
leads to a contradiction (one obtains $\alpha<\alpha$).

\medskip
\noindent For $\alpha=|A|$ and $\beta=|B|$,
\[
\alpha\leq\beta\ \iff
\exists\, f:A\longrightarrow B\text{ injective}.
\]
The strict inequality is given by
\[
\alpha<\beta\quad\iff\quad
\begin{cases}
\text{there exists an injection }f:A\longrightarrow B,\\
\text{there is no bijection }A\longrightarrow B.
\end{cases}
\]

\noindent The definition is correct.
\endgroup
; no preamble or page break.
% EDITORIAL NOTES:
% - This final manuscript page continues the discussion of cardinality
%   begun on Seminar 3, Page 9.
% - For the cardinal 1 the handwritten set looks like {{empty set}}, rather
%   than the usual {empty set}. Either is a singleton and has cardinal 1.
% - The handwritten cardinal sum appears to use A_i UNION {i}, not
%   A_i TIMES {i}. That literal expression does not in general produce a
%   disjoint union, and its cardinality need not be the intended sum.
%   The standard tagged-disjoint-union construction would instead be
%     | \bigcup_{i\in I} (A_i \times \{i\}) |.
%   We retain the apparent handwriting in the displayed expression.
% - The manuscript states that these definitions are correct without a
%   proof of independence of the representing sets A_i, A and B.
% - "Cantor's paradox" is transcribed as a note about the putative set
%   of all cardinals; its proper-class distinction is not supplied in
%   the manuscript and is not inserted into the text.
% - The last comparison of cardinalities uses injective/bijective maps
%   without addressing any set-theoretic axioms. It is left as written.

\subsection*{Seminar 3 - Page 10}
\begingroup
\setlength{\abovedisplayskip}{5pt}
\setlength{\belowdisplayskip}{5pt}
\setlength{\abovedisplayshortskip}{3pt}
\setlength{\belowdisplayshortskip}{3pt}

\noindent The first cardinal numbers are written as
\[
\begin{aligned}
0 &:= \operatorname{card}(\varnothing),\\
1 &:= \operatorname{card}\bigl(\{\{\varnothing\}\}\bigr),\\
2 &:= \operatorname{card}\bigl(\{\{\varnothing\},\{\{\varnothing\}\}\}\bigr).
\end{aligned}
\]

\medskip
\noindent\underline{Operations with cardinal numbers.}

Let $\alpha=|A|$. If $(\alpha_i)_{i\in I}$ is a family of cardinal
numbers, with $\alpha_i=|A_i|$, define
\[
\sum_{i\in I}\alpha_i\ \overset{\mathrm{def}}{=}
\left|\bigcup_{i\in I}\bigl(A_i\cup\{i\}\bigr)\right|.
\]

\noindent For $\alpha=|A|$ and $\beta=|B|$, define
\[
\alpha\cdot\beta\ \overset{\mathrm{def}}{=}\ |A\times B|,
\qquad
\alpha^{\beta}\ \overset{\mathrm{def}}{=}
\bigl|\operatorname{Hom}(B,A)\bigr|.
\]

\noindent The definitions are well-defined.

\medskip
\noindent\underline{Cantor's paradox.}
The collection of all cardinal numbers, regarded as a set,
leads to a contradiction (one obtains $\alpha<\alpha$).

\medskip
\noindent For $\alpha=|A|$ and $\beta=|B|$,
\[
\alpha\leq\beta\ \iff
\exists\, f:A\longrightarrow B\text{ injective}.
\]
The strict inequality is given by
\[
\alpha<\beta\quad\iff\quad
\begin{cases}
\text{there exists an injection }f:A\longrightarrow B,\\
\text{there is no bijection }A\longrightarrow B.
\end{cases}
\]

\noindent The definition is correct.
\endgroup
; no preamble or page break.
% EDITORIAL NOTES:
% - This final manuscript page continues the discussion of cardinality
%   begun on Seminar 3, Page 9.
% - For the cardinal 1 the handwritten set looks like {{empty set}}, rather
%   than the usual {empty set}. Either is a singleton and has cardinal 1.
% - The handwritten cardinal sum appears to use A_i UNION {i}, not
%   A_i TIMES {i}. That literal expression does not in general produce a
%   disjoint union, and its cardinality need not be the intended sum.
%   The standard tagged-disjoint-union construction would instead be
%     | \bigcup_{i\in I} (A_i \times \{i\}) |.
%   We retain the apparent handwriting in the displayed expression.
% - The manuscript states that these definitions are correct without a
%   proof of independence of the representing sets A_i, A and B.
% - "Cantor's paradox" is transcribed as a note about the putative set
%   of all cardinals; its proper-class distinction is not supplied in
%   the manuscript and is not inserted into the text.
% - The last comparison of cardinalities uses injective/bijective maps
%   without addressing any set-theoretic axioms. It is left as written.

\subsection*{Seminar 3 - Page 10}
\begingroup
\setlength{\abovedisplayskip}{5pt}
\setlength{\belowdisplayskip}{5pt}
\setlength{\abovedisplayshortskip}{3pt}
\setlength{\belowdisplayshortskip}{3pt}

\noindent The first cardinal numbers are written as
\[
\begin{aligned}
0 &:= \operatorname{card}(\varnothing),\\
1 &:= \operatorname{card}\bigl(\{\{\varnothing\}\}\bigr),\\
2 &:= \operatorname{card}\bigl(\{\{\varnothing\},\{\{\varnothing\}\}\}\bigr).
\end{aligned}
\]

\medskip
\noindent\underline{Operations with cardinal numbers.}

Let $\alpha=|A|$. If $(\alpha_i)_{i\in I}$ is a family of cardinal
numbers, with $\alpha_i=|A_i|$, define
\[
\sum_{i\in I}\alpha_i\ \overset{\mathrm{def}}{=}
\left|\bigcup_{i\in I}\bigl(A_i\cup\{i\}\bigr)\right|.
\]

\noindent For $\alpha=|A|$ and $\beta=|B|$, define
\[
\alpha\cdot\beta\ \overset{\mathrm{def}}{=}\ |A\times B|,
\qquad
\alpha^{\beta}\ \overset{\mathrm{def}}{=}
\bigl|\operatorname{Hom}(B,A)\bigr|.
\]

\noindent The definitions are well-defined.

\medskip
\noindent\underline{Cantor's paradox.}
The collection of all cardinal numbers, regarded as a set,
leads to a contradiction (one obtains $\alpha<\alpha$).

\medskip
\noindent For $\alpha=|A|$ and $\beta=|B|$,
\[
\alpha\leq\beta\ \iff
\exists\, f:A\longrightarrow B\text{ injective}.
\]
The strict inequality is given by
\[
\alpha<\beta\quad\iff\quad
\begin{cases}
\text{there exists an injection }f:A\longrightarrow B,\\
\text{there is no bijection }A\longrightarrow B.
\end{cases}
\]

\noindent The definition is correct.
\endgroup
; no preamble or page break.
% EDITORIAL NOTES:
% - This final manuscript page continues the discussion of cardinality
%   begun on Seminar 3, Page 9.
% - For the cardinal 1 the handwritten set looks like {{empty set}}, rather
%   than the usual {empty set}. Either is a singleton and has cardinal 1.
% - The handwritten cardinal sum appears to use A_i UNION {i}, not
%   A_i TIMES {i}. That literal expression does not in general produce a
%   disjoint union, and its cardinality need not be the intended sum.
%   The standard tagged-disjoint-union construction would instead be
%     | \bigcup_{i\in I} (A_i \times \{i\}) |.
%   We retain the apparent handwriting in the displayed expression.
% - The manuscript states that these definitions are correct without a
%   proof of independence of the representing sets A_i, A and B.
% - "Cantor's paradox" is transcribed as a note about the putative set
%   of all cardinals; its proper-class distinction is not supplied in
%   the manuscript and is not inserted into the text.
% - The last comparison of cardinalities uses injective/bijective maps
%   without addressing any set-theoretic axioms. It is left as written.

\subsection*{Seminar 3 - Page 10}
\begingroup
\setlength{\abovedisplayskip}{5pt}
\setlength{\belowdisplayskip}{5pt}
\setlength{\abovedisplayshortskip}{3pt}
\setlength{\belowdisplayshortskip}{3pt}

\noindent The first cardinal numbers are written as
\[
\begin{aligned}
0 &:= \operatorname{card}(\varnothing),\\
1 &:= \operatorname{card}\bigl(\{\{\varnothing\}\}\bigr),\\
2 &:= \operatorname{card}\bigl(\{\{\varnothing\},\{\{\varnothing\}\}\}\bigr).
\end{aligned}
\]

\medskip
\noindent\underline{Operations with cardinal numbers.}

Let $\alpha=|A|$. If $(\alpha_i)_{i\in I}$ is a family of cardinal
numbers, with $\alpha_i=|A_i|$, define
\[
\sum_{i\in I}\alpha_i\ \overset{\mathrm{def}}{=}
\left|\bigcup_{i\in I}\bigl(A_i\cup\{i\}\bigr)\right|.
\]

\noindent For $\alpha=|A|$ and $\beta=|B|$, define
\[
\alpha\cdot\beta\ \overset{\mathrm{def}}{=}\ |A\times B|,
\qquad
\alpha^{\beta}\ \overset{\mathrm{def}}{=}
\bigl|\operatorname{Hom}(B,A)\bigr|.
\]

\noindent The definitions are well-defined.

\medskip
\noindent\underline{Cantor's paradox.}
The collection of all cardinal numbers, regarded as a set,
leads to a contradiction (one obtains $\alpha<\alpha$).

\medskip
\noindent For $\alpha=|A|$ and $\beta=|B|$,
\[
\alpha\leq\beta\ \iff
\exists\, f:A\longrightarrow B\text{ injective}.
\]
The strict inequality is given by
\[
\alpha<\beta\quad\iff\quad
\begin{cases}
\text{there exists an injection }f:A\longrightarrow B,\\
\text{there is no bijection }A\longrightarrow B.
\end{cases}
\]

\noindent The definition is correct.
\endgroup
